\begin{chaptersummary}
  \item Composition reads schema files and never calls a service. That is why
        a subgraph does not have to be running to be composed, why a document
        no service exports can still be composed against yours, and why the
        \code{routing\_url} in the input file is written into the output rather
        than fetched from.
  \item The first thing that stopped this graph was not federation. Hot
        Chocolate applies \code{@cost} to every field and \code{@listSize} to a
        paged one with no configuration at all, and the composer holds its own
        definitions of both, so six weights and one argument were rejected
        before anything about a seam was reached.
        \code{ApplyCostDefaults = false} is the whole fix, and what it removes
        was a description of a paging limit rather than the limit itself.
  \item A router execution config is not a supergraph schema; it contains one.
        Its shape is defined by a protobuf in the Cosmo source and by no
        documentation, and it is worth opening once. Inside \code{engineConfig}
        are two schemas rather than one: the client's, with every federation
        directive gone, and each subgraph's own document with them all still
        there, next to the keys compiled into a lookup table. The top level
        carries the version of the composition library that produced the file,
        which is the only record of that anywhere.
  \item \code{Query.node} needed \code{@shareable} and there was no method to
        put it on. The attribute on the \code{Query} class works and declares
        every root field shareable to do it, which switches off the check that
        stops another subgraph claiming \code{sessionById}. A type interceptor
        is thirty-three lines and reaches exactly the two generated fields,
        and those lines are what the narrower promise costs.
  \item Composition is two passes with different powers. Merging decides
        whether the documents contradict each other and cannot be skipped.
        Resolvability walks out from each root field asking whether a plan
        exists for every field it reaches, prints the path it got stuck on, and
        can be switched off with a flag. A graph that composes has no
        contradiction the composer can see, which is not the same as being
        right: \code{@shareable} on a field only one service owns composes
        cleanly and is a lie.
\end{chaptersummary}
